\section{Topic Introduction}

\subsection{Statement of the Problem}

For much of the twentieth century, space exploration was the exclusive domain of sovereign states. The “space age” was defined by government-led programmes, funded by national treasuries, staffed by government scientists and astronauts, and conducted in pursuit of geopolitical prestige and scientific discovery.\footnote{Belfer Center for Science and International Affairs. (n.d.). \emph{The new space race}. Harvard Kennedy School. \url{https://www.belfercenter.org/research-analysis/new-space-race}} The legal framework that emerged during this era, most notably the 1967 Outer Space Treaty, reflected these assumptions: space was to be used for the benefit and in the interests of all countries, free from national appropriation.\footnote{United Nations Office for Outer Space Affairs. (n.d.). \emph{The Outer Space Treaty}. \url{https://www.unoosa.org/oosa/en/ourwork/spacelaw/treaties/introouterspacetreaty.html}}

In the twenty-first century, however, a fundamental transformation has taken place. Private companies now routinely launch satellites, resupply the International Space Station, and are developing the capacity to carry humans to the Moon, Mars, and beyond. The emergence of what is often called “NewSpace”, commercially driven space ventures backed by private capital, has reframed outer space not merely as a frontier of exploration but as a potential arena of economic activity, resource extraction, and permanent human settlement.\footnote{Belfer Center for Science and International Affairs. (n.d.). \emph{The new space race}. Harvard Kennedy School. \url{https://www.belfercenter.org/research-analysis/new-space-race}}

This shift raises profound questions for legislators. The central tension can be stated as follows: how should the United States encourage private investment, innovation, and economic and strategic leadership in space whilst honouring the international-law principles that outer space is the “province of all mankind,” not subject to national appropriation, and must be used for the benefit of all countries?\footnote{United Nations Office for Outer Space Affairs. (n.d.). \emph{The Outer Space Treaty}. \url{https://www.unoosa.org/oosa/en/ourwork/spacelaw/treaties/introouterspacetreaty.html}}

\subsection{Key Flashpoints}

The committee should be aware of the following principal areas of contention:

\begin{enumerate}
  \item \textbf{Space Resource and Asteroid Mining Property Rights:} Can private entities own resources extracted from celestial bodies? How does this interact with the non-appropriation principle?
  \item \textbf{Orbital Debris and Space Traffic Management:} As the number of private satellites and objects in orbit proliferates, who is responsible for managing orbital traffic and mitigating collision risks?
  \item \textbf{Mission Authorisation and Regulatory Gaps:} Many novel private-sector space activities (in-space servicing, manufacturing, debris removal) lack clear governance under existing U.S. regulation.
  \item \textbf{Equity Between Spacefaring and Non-Spacefaring Nations:} Does a national legislative approach to space resources risk creating a “first-come, first-served” regime that disadvantages developing nations?
  \item \textbf{Permanent Human Settlement (“Colonisation”):} What legal, ethical, and governance frameworks should apply to the prospect of human communities on celestial bodies?
\end{enumerate}

\subsection{The U.S. Legislative Context}

The United States has been at the forefront of legislative efforts to open outer space to commercial activity.\footnote{von der Dunk, F. (2015). \emph{The US Space Launch Competitiveness Act of 2015}. Space, Cyber, and Telecommunications Law Program Faculty Publications, University of Nebraska. \url{https://digitalcommons.unl.edu/spacelaw/98/}} Through a series of statutes, from the Commercial Space Launch Act of 1984 to the SPACE Act of 2015 and subsequent bills, Congress has progressively expanded the legal basis for private-sector operations in outer space, including the recognition of property rights in extracted space resources.\footnote{Congressional Research Service. (n.d.). \emph{Commercial space launch and reentry regulations: Overview and select issues} (R48582). Congress.gov. \url{https://www.congress.gov/crs-product/R48582}} These domestic measures interact, sometimes uneasily, with the international legal framework that governs outer space.\footnote{Hardman, J. (2023). For all mankind? The United States' race to redefine space law. \emph{UCLA Law Review}. \url{https://uclawreview.org/2023/03/23/united-states-race-to-redefine-space-law/}}
